\section{Related Work}
\label{sec:related}

\paragraph{RL-trained LLM agents.}
RL fine-tuning endows LLMs with procedural capabilities SFT does not produce~\cite{zhou2025mem1,zhou2024webarena,liu2024agentbench,yao2023react,zeng2024agenttuning}. We investigate why these RL-acquired strategies resist distillation.

\paragraph{Knowledge distillation for LLMs.}
Standard methods~\cite{hsieh2023distilling,gu2024minillm,jiang2023lion} target single-turn transfer; policy distillation and imitation learning~\cite{rusu2016policy,czarnecki2019distilling,ross2011dagger} handle procedural knowledge in low-dimensional action spaces. Our setting differs: the teacher's competence arises from a \emph{multi-turn procedural policy} in the high-dimensional token space of LLMs.

\paragraph{SFT--RL alignment gaps and weight-space analysis.}
The ``alignment tax''~\cite{askell2021alignment} documents RLHF degrading pre-training capabilities. \citet{jin2025rl} show SFT induces ``hard alignment'' while RL produces ``soft alignment''; \citet{aghajanyan2021intrinsic} demonstrate fine-tuning operates in low intrinsic dimension. PiSSA~\cite{meng2024pissa} initializes LoRA from principal singular components; our SVD alignment ratio measures the overlap of \emph{fine-tuning deltas} with the base model's principal subspace. We extend these by revealing quantitative gaps ($2$--$6\times$ for RL vs.\ $1.5\times$ for SFT) not captured by prior metrics. Extended related work on multi-objective RL and preference learning is in Appendix~\ref{app:extended_related}.

\paragraph{The evaluation blind spot in agentic distillation.}
Recent work increasingly distills tool-using capabilities from large models into smaller ones, but evaluation practices vary widely in whether they test actual tool execution. Several works evaluate only format correctness: \citet{jhandi2025slm} use ChatGPT-as-judge to score tool-call syntax without executing calls; \citet{yin2025magnet} evaluate against BFCL format matching, which measures argument-level accuracy but not end-to-end execution; AgentKD~\cite{agentkd2025} trains on virtual tool-use tasks with pre-defined return values. These evaluations implicitly assume that format accuracy entails execution capability. Our results show this assumption can fail catastrophically: Jaccard $= 0.95$ (near-perfect format) coexists with $\text{em\_partial} = 0.025$ (near-zero task performance). Works that provide real tool access at test time~\cite{kang2025agent,liu2025structured} produce more robust evaluations but still lack a with/without tool-backend ablation that would reveal whether the student genuinely depends on tool execution. To our knowledge, no prior agentic distillation work reports an oracle retrieval ablation that directly compares student performance with and without the tool backend on the same inputs---the diagnostic that directly quantifies the execution gap. We propose this ablation as a standard component of agentic distillation evaluation.
